\chapter{Das Siegelbandlogo}

Die verschiedenen Farbvarianten des Siegelbandlogos in \tubslatex werden
vom Paket \lstinline{tubslogo} zur Verfügung gestellt.
Die folgende Beschreibung bezieht sich auf den Funktionsumfang des Paketes.
Das Paket ist in allen verfügbaren Klassen bereits korrekt eingebunden.

Zur Verfügung gestellt werden eine Variante im rgb-Farbmodell, eine im
cmyk-Farbmodell, eine mit der Sonderfarbe HKS 15 oder (für Ausnahmefälle)
eine monochromen Variante.
Die Auwahl erfolgt dabei über Paketoptionen
(siehe Abschnitt \ref{options:color}).

\begin{figure}[!ht]
\begin{minipage}{0.5\textwidth}
  \centering
  \includegraphics[width=\tubslogoBaseWidth]{TUBraunschweig-rgb}
  \subcaption{Logo in Standardgröße (RGB)}
\end{minipage}
\begin{minipage}{0.5\textwidth}
  \centering
  \includegraphics[width=\tubslogoBaseWidth]{TUBraunschweig_4C}
  \subcaption{Logo in Standardgröße (CMYK)}
\end{minipage}

\begin{minipage}{0.5\textwidth}
  \centering
  \includegraphics[width=\tubslogoBaseWidth]{TUBraunschweig_SC}
  \subcaption{Logo in Standardgröße (HKS 15)}
\end{minipage}
\begin{minipage}{0.5\textwidth}
  \centering
  \includegraphics[width=\tubslogoBaseWidth]{TUBraunschweig_B}
  \subcaption{Logo in Standardgröße (schwarz)}
\end{minipage}
\caption{Verfügbare Farbversionen der Logos}
\end{figure}

\begin{Declaration}
  \Macro{tubslogo}\OParameter{Skalierung}\\
  \Macro{tubslogoAbs}\Parameter{Breite}
\end{Declaration}


Es werden 2 Befehle zum Einfügen des Logos in Dokumente bereit gestellt.
Mit dem Befehl \Macro{tubslogo} wird ein Logo in
standardkonformer Größe der verwendeten Papiergröße entsprechend skaliert
eingefügt. Die verwendete Papiergröße kann dabei als Paketoption übergeben
werden und ist standardmäßig auf \lstinline{a4paper} voreingestellt
(siehe Abschnitt \ref{options:papersize}).
Als optionales Argument kann auch ein frei gewählter Skalierungsfaktor
vorgegeben werden (siehe Abschnitt \ref{cmd:tubslogo}).

Mit \Macro{tubslogoAbs} kann über den Parameter \PName{Breite} dagegen
zusätzlich eine individuell gewählte absolute Breite angegeben werden.


\section{Paket-Optionen}

\subsection{Farbmodell}\label{options:color}

\begin{description}
  \classoptionitem[rgb]
    Verwendung der rgb-Druckfarben. Dies ist die Standardeinstellung.
  \classoptionitem[cmyk]
    Verwendung der cmyk-Farben
  \classoptionitem[hks]
    Verwendung der Spezialfarbe HKS 15.
  \classoptionitem[mono]
    Verwendung einer schwarz-weiß-Version des Logos.
\end{description}

\subsection{Papiergröße}\label{options:papersize}

\begin{description}
  \classoptionitem[a6paper]
    Skaliert das Logo passend für A6-Dokumente auf 60\% der Standardgröße.
  \classoptionitem[a5paper]
    Skaliert das Logo passend für A5-Dokumente auf 70\% der Standardgröße.
  \classoptionitem[a4paper]
    Skaliert das Logo passend für A4-Dokumente auf Standardgröße.
  \classoptionitem[a3paper]
    Skaliert das Logo passend für A3-Dokumente auf 140\% der Standardgröße.
  \classoptionitem[a2paper]
    Skaliert das Logo passend für A2-Dokumente auf 200\% der Standardgröße.
  \classoptionitem[a1paper]
    Skaliert das Logo passend für A1-Dokumente auf 280\% der Standardgröße.
  \classoptionitem[a0paper]
    Skaliert das Logo passend für A0-Dokumente auf 400\% der Standardgröße.
\end{description}


\clearpage
\section{Befehle}

\begin{description}\label{cmd:tubslogo}
  \item[\color{tuRed}\mdseries\ttfamily \textbackslash tubslogo%
    {[\textcolor{tuGreenDark}{\sffamily\itshape $<$Breite rel.$>$}]}]
    Erzeugt ein Logo, dass entsprechend des durch Angabe der Papiergröße
    gewählten Skalierungsfaktors skaliert wird.

    Wird das optionale Argument {\sffamily\itshape $<$Breite rel.$>$}
    verwendet, so kann ein freier Skalierungsfaktor ausgehend von der
    Standardgröße für DINA4-Dokumente ($70$mm$\times 26$mm) übergeben werden.

    \lstinline!\tubslogo[0.6]! erzeugt beispielsweise ein auf 60\% der 
    Standardgröße für DINA4-Dokumente skaliertes Logo.

  \item[\color{tuRed}\mdseries\ttfamily \textbackslash tubslogoAbs%
    \{\textcolor{tuGreenDark}{Breite abs.}\}]
    Dieser Befehl erlaubt die Erzeugung eines Logos mit absolut vorgegebener
    Breite.\footnote{Generell ist die Verwendung von standardkonformen Größen
    einer individuellen Skalierung vorzuziehen!}

    \lstinline!\tubslogoAbs{5cm}! erzeugt beispielsweise ein 5cm breites Logo.
\end{description}


\section{Längen}
  \begin{description}
    \item[\mdseries\ttfamily \textbackslash tubslogoBaseWidth]\hfill\\
      Standardbreite des Logos (70mm).
    \item[\mdseries\ttfamily \textbackslash tubslogoBaseHeight]\hfill\\
      Standardhöhe des Logos (26mm).
    \item[\mdseries\ttfamily \textbackslash tubslogoWidth]\hfill\\
      Breite des automatisch skalierten Logos. Wird automatisch berechnet und
      kann zur dynamischen Positionierung verwendet werden.
    \item[\mdseries\ttfamily \textbackslash tubslogoHeight]\hfill\\
      Höhe des automatisch skalierten Logos. Siehe \lstinline{tubslogoWidth}.
  \end{description}
